\documentclass[14pt,a4paper,oneside]{extbook}
\usepackage[utf8]{vietnam}
\usepackage{amsmath,amssymb,mathrsfs,fancyhdr,enumerate,multirow,makecell}
\usepackage{tkz-euclide,tikz-3dplot,tikz,tkz-tab}
\usetikzlibrary{arrows.meta,calc,intersections,angles,quotes,backgrounds,shapes.geometric,patterns,through,shadows,positioning}
\usepackage{yhmath}
\usepackage[top=0.7cm, bottom=0.7cm, left=1.0cm, right=1.0cm, headheight=16pt, headsep=3mm, footskip=5mm]{geometry}
\usepackage[hidelinks,unicode]{hyperref}
\usepackage{qrcode}
\usepackage[many]{tcolorbox}
\usepackage{fontawesome}

\pagestyle{fancy}
\fancyhf{}
\fancyhead[L]{\footnotesize\bfseries\color{teal!85!black}\faBook~SỔ TAY KIẾN THỨC TOÁN 5}
\fancyhead[R]{\footnotesize\bfseries\color{blue!85!black}Lớp Toán Thầy Lĩnh -- loptoanthaylinh.vn}
\fancyfoot[C]{\footnotesize\bfseries\color{gray!80!black}-- Sổ tay kiến thức: Bài 54. Thực hành tính toán, ước lượng thể tích --}
\renewcommand{\headrulewidth}{0.4pt}

\begin{document}

% --- BANNER TIÊU ĐỀ BÀI HỌC ---
\noindent
\begin{tcolorbox}[
  enhanced,
  colback=teal!8!white,
  colframe=teal!80!black,
  arc=3mm,
  boxrule=1.1pt,
  width=\linewidth,
  halign=center,
  top=2.2mm, bottom=2.2mm,
  left=3mm, right=3mm,
  drop fuzzy shadow
]
  {\large\bfseries\color{teal!95!black}\faBookmark\ BÀI 54. THỰC HÀNH TÍNH TOÁN, ƯỚC LƯỢNG THỂ TÍCH}\\[2pt]
  {\scriptsize\color{black!75}\itshape Đo đạc kích thước thực tế, tính dung tích đồ vật và ước lượng thể tích cuộc sống trên 1 trang A4}
\end{tcolorbox}

\vspace{3.2mm}
% --- KHỐI 1: CHỌN ĐƠN VỊ THỂ TÍCH THỰC TẾ ---
\noindent
\begin{minipage}[t]{0.54\textwidth}
  \begin{tcolorbox}[
    enhanced,
    title={\footnotesize\bfseries\color{white}\faEye\ 1. Chọn đơn vị thể tích phù hợp},
    colback=blue!2!white,
    colframe=blue!75!black,
    arc=2.5mm, boxrule=0.8pt, top=1.8mm, bottom=1.8mm, left=2.2mm, right=2.2mm
  ]
    \footnotesize
    $\bullet$ \textbf{Xăng-ti-mét khối ($\text{cm}^3$ / ml):} Thể tích chén nước, hộp sữa chua, hộp thuốc tiêm, viên xúc xắc.\\
    $\bullet$ \textbf{Đề-xi-mét khối ($\text{dm}^3$ / Lít):} Chai nước suối, xô nước, nồi cơm điện, ba lô học sinh.\\
    $\bullet$ \textbf{Mét khối ($\text{m}^3$):} Bể chứa nước gia đình, thể tích phòng học, thùng xe tải, bể bơi.\\
    \textit{VD:} Một chai nước ngọt cỡ vừa thường có dung tích $1,5\text{ lít} = 1,5\text{ dm}^3$.
  \end{tcolorbox}
\end{minipage}%
\hfill
\begin{minipage}[t]{0.44\textwidth}
  \begin{tcolorbox}[
    enhanced,
    title={\footnotesize\bfseries\color{white}\faRulerCombined\ 2. Kỹ năng đo đạc vật thể},
    colback=teal!2!white,
    colframe=teal!75!black,
    arc=2.5mm, boxrule=0.8pt, top=1.8mm, bottom=1.8mm, left=2mm, right=2mm
  ]
    \footnotesize
    $\bullet$ \textbf{Đo bên ngoài}: Cho ra thể tích toàn bộ vật chiếm chỗ trong không gian.\\
    $\bullet$ \textbf{Đo bên trong (lòng hộp)}: Cho ra \textbf{dung tích} (khả năng chứa đựng chất lỏng, hàng hóa bên trong).
  \end{tcolorbox}
\end{minipage}

\vspace{3mm}

% --- KHỐI 2: DẠNG TOÁN XẾP HỘP NHỎ VÀO THÙNG LỚN ---
\noindent
\begin{tcolorbox}[
  enhanced,
  title={\footnotesize\bfseries\color{white}\faTh\ 3. Bài toán xếp các hộp nhỏ vào thùng carton lớn},
  colback=purple!2!white,
  colframe=purple!70!black,
  arc=2.5mm, boxrule=0.8pt, top=1.8mm, bottom=1.8mm, left=2.5mm, right=2.5mm
]
  \footnotesize
  Một thùng carton kích thước $60\text{ cm} \times 40\text{ cm} \times 30\text{ cm}$. Người ta xếp các hộp chè hình lập phương cạnh $10\text{ cm}$ vào thùng. Hỏi xếp được nhiều nhất bao nhiêu hộp chè?\\
  $\triangleright$ Số hộp xếp theo chiều dài: $60 : 10 = 6\text{ (hộp)}$.\\
  $\triangleright$ Số hộp xếp theo chiều rộng: $40 : 10 = 4\text{ (hộp)}$. $\rightarrow$ Mỗi lớp xếp được: $6 \times 4 = 24\text{ (hộp)}$.\\
  $\triangleright$ Số lớp xếp theo chiều cao: $30 : 10 = 3\text{ (lớp)}$.\\
  $\triangleright$ \textbf{Tổng số hộp chè xếp được}: $24 \times 3 = \mathbf{72\text{ hộp}}$ (hoặc $V_{\text{thùng}} : V_{\text{hộp}} = 72\,000 : 1\,000 = 72$).
\end{tcolorbox}

\vspace{3mm}

% --- KHỐI 3: SƠ ĐỒ TƯ DUY ---
\noindent
\begin{tcolorbox}[
  enhanced,
  title={\footnotesize\bfseries\color{white}\faSitemap\ Sơ đồ tư duy: Ước lượng \& Đóng gói hình khối},
  colback=yellow!2!white,
  colframe=teal!85!black,
  arc=3mm, boxrule=0.9pt, top=1.8mm, bottom=1.8mm, left=1.5mm, right=1.5mm,
  halign=center
]
  \begin{tikzpicture}[
    scale=0.92,
    every node/.style={transform shape},
    b1/.style={rectangle, rounded corners=4pt, fill=blue!15, draw=blue!70!black, text=blue!90!black, font=\bfseries\scriptsize, inner sep=4pt, align=center},
    b2/.style={rectangle, rounded corners=4pt, fill=teal!15, draw=teal!70!black, text=teal!90!black, font=\bfseries\scriptsize, inner sep=4pt, align=center},
    b3/.style={rectangle, rounded corners=4pt, fill=green!15, draw=green!70!black, text=green!90!black, font=\bfseries\scriptsize, inner sep=4pt, align=center}
  ]
    \node[b1] (n1) at (-5.8, 0) {1. Thùng lớn: $A \times B \times C$\\Đổi cùng đơn vị đo};
    \node[b2] (n2) at (0, 0) {2. Hộp nhỏ: $a \times b \times c$\\Chia theo từng chiều};
    \node[b3] (n3) at (5.8, 0) {3. Số hộp tối đa:\\$(A:a) \times (B:b) \times (C:c)$};

    \draw[teal!85!black, thick, -{Stealth[length=2.5mm]}] (n1) -- (n2);
    \draw[teal!85!black, thick, -{Stealth[length=2.5mm]}] (n2) -- (n3);
  \end{tikzpicture}
\end{tcolorbox}

\vspace{3mm}

% --- KHỐI 4: LƯU Ý VỀ ĐỘ DÀY THÀNH BỂ & BẪY TRẮC NGHIỆM ---
\noindent
\begin{minipage}[t]{0.54\textwidth}
  \begin{tcolorbox}[
    enhanced,
    title={\footnotesize\bfseries\color{white}\faShield\ Bể có độ dày của thành và đáy},
    colback=green!3!white,
    colframe=green!60!black,
    arc=2.5mm, boxrule=0.8pt, top=1.8mm, bottom=1.8mm, left=2.2mm, right=2.2mm
  ]
    \footnotesize
    Nếu bể có thành dày $d$ và đáy dày $d$:\\
    $\bullet$ Chiều dài trong lòng: $a' = a - 2d$.\\
    $\bullet$ Chiều rộng trong lòng: $b' = b - 2d$.\\
    $\bullet$ Chiều cao trong lòng (không nắp): $c' = c - d$.\\
    $\triangleright$ Dung tích nước chứa được phải tính theo kích thước trong lòng bể!
  \end{tcolorbox}
\end{minipage}%
\hfill
\begin{minipage}[t]{0.44\textwidth}
  \begin{tcolorbox}[
    enhanced,
    title={\footnotesize\bfseries\color{white}\faExclamationTriangle\ Mẹo tránh bẫy trắc nghiệm},
    colback=red!2!white,
    colframe=red!70!black,
    arc=2.5mm, boxrule=0.8pt, top=1.8mm, bottom=1.8mm, left=2.2mm, right=2.2mm
  ]
    \footnotesize
    $\bullet$ Không thể chỉ lấy thể tích thùng chia thể tích hộp nếu kích thước các cạnh không chia hết (sẽ bị khoảng trống không xếp được)!\\
    $\bullet$ Phải chia theo từng kích thước dài, rộng, cao.
  \end{tcolorbox}
\end{minipage}

\end{document}
